\documentclass[aspectratio=169]{beamer}

\usetheme{Madrid}
\usecolortheme{default}
\usepackage{tikz}
\usepackage{amsmath}
\usetikzlibrary{shapes.geometric, arrows.meta, positioning, fit, calc}

% Define TikZ styles
\tikzstyle{block} = [rectangle, rounded corners, minimum width=2.5cm, minimum height=1cm,text centered, draw=black, fill=blue!10]
\tikzstyle{arrow} = [thick,->,>=stealth]
\tikzstyle{process} = [rectangle, draw=black, fill=orange!10, text centered, rounded corners, minimum height=0.8cm, inner sep=4pt]

\title[Soft Routes Integration]{Theoretical Models for Integrating Soft Routes\\ into the Ride-Sharing Cost Function}
\author[Intern Name]{Intern Name}
\institute[University/Lab]{University/Lab}
\date{\today}

\begin{document}

\begin{frame}
  \titlepage
\end{frame}

\begin{frame}{Motivation: Moving Beyond Heuristics}
  \textbf{The Problem with Hardcoded Rules:}
  \begin{itemize}
      \item Forcing vehicles to park specifically on predefined trunks (or bypassing feasibility checks if they share a trunk) is an \textbf{engineering heuristic}.
      \item It artificially frames the simulation results, leading to unrealistic traffic distributions and removing the natural flexibility of the ride-sharing algorithm.
  \end{itemize}
  
  \vspace{0.3cm}
  \textbf{The Academic Objective:}
  \begin{itemize}
      \item We must build a \textbf{mathematical model} where soft routes (high-demand corridors extracted via topology scanning) are integrated organically into the objective optimization engine.
      \item By bringing soft routes into the \textbf{Cost Function}, the fleet naturally gravitates towards these corridors because the math makes it the globally optimal choice.
  \end{itemize}
\end{frame}

\begin{frame}{The Current Baseline: Version 1 Base Cost Function}
  Before integrating the soft routes, we must understand the baseline optimization objective. The simulation evaluates any passenger subset $S$ assigned to a vehicle $v$ by minimizing the total penalty cost:
  
  \vspace{0.3cm}
  \begin{block}{Base Objective Cost Function}
  \[ C_{base}(v, S) = \underbrace{\gamma \cdot \Delta_{new}}_{\text{Detour Cost}} + \underbrace{\alpha \sum_{p \in S} d_{walk}(p)}_{\text{Walking Penalty}} + \underbrace{\beta \sum_{p \in S} \delta_{delay}(p)}_{\text{Trip Delay}} + \underbrace{\mu(Cap_v - |S|)}_{\text{Unused Capacity}} \]
  \end{block}
  
  \vspace{0.2cm}
  \begin{itemize}
      \item The algorithm strictly minimizes distance, walking discomfort, and wait times.
      \item \textbf{Limitation:} It is completely "blind" to future demand patterns. A 500m detour into a high-demand downtown corridor costs exactly the same as a 500m detour into a dead-zone.
  \end{itemize}
\end{frame}

\begin{frame}{Approach 1: The Corridor Adherence Reward (Path-Based)}
  \textbf{Concept:} Introduce a dynamic reward term that subsidizes the base cost if the proposed subset's path travels along the edges of an active soft route.
  
  \vspace{0.3cm}
  \begin{block}{Mathematical Formulation}
  Let $\Omega(R)$ represent the fraction of the proposed vehicle route $R$ that overlaps with the set of active soft route edges $E_{soft}$.
  \[ C_{new} = C_{base}(v, S) \textbf{ } \mathbf{-} \textbf{ } \lambda \cdot \Omega(R) \]
  Where $\lambda$ is a tunable reward coefficient.
  \end{block}
  
  \vspace{0.3cm}
  \textbf{Algorithmic Logic:}
  \begin{itemize}
      \item The optimizer naturally prefers to pick up passengers whose trips run parallel to the soft routes.
      \item This effectively "subsidizes" the detour cost: a longer detour that drops a passenger along a trunk might become mathematically cheaper than a short detour into a dead-zone.
  \end{itemize}
\end{frame}

\begin{frame}{Approach 2: Spatially-Weighted Detour Cost (Edge-Based)}
  \textbf{Concept:} Instead of applying a flat penalty for detour distance, we alter how the vehicle "perceives" distance based on its geographical location.
  
  \vspace{0.3cm}
  \begin{block}{Mathematical Formulation}
  We replace the flat detour term $\gamma \cdot \Delta_{new}$ with a discrete, edge-by-edge summation:
  \[ DetourCost = \gamma \cdot \sum_{e \in Route} \Big( length(e) \cdot W(e) \Big) \]
  \[ 
    W(e) = 
    \begin{cases} 
      \omega_{soft} & \text{if } e \in E_{soft} \quad (\text{e.g., } 0.5) \\
      1.0 & \text{otherwise} 
    \end{cases}
  \]
  \end{block}
  
  \vspace{0.2cm}
  \textbf{Algorithmic Logic:}
  \begin{itemize}
      \item Driving 500 meters along a soft route is calculated by the optimizer as only 250 meters.
      \item The algorithm will aggressively route empty vehicles towards soft routes to pick up passengers, because traversing that space incurs less mathematical penalty.
  \end{itemize}
\end{frame}

\begin{frame}{Approach 3: Terminal Opportunity Cost (Destination-Based)}
  \textbf{Concept:} The system ignores \textit{how} the vehicle travels, but heavily scrutinizes \textit{where it ends up}. Dropping off a passenger in a dead-zone incurs an opportunity cost.
  
  \vspace{0.3cm}
  \begin{block}{Mathematical Formulation}
  Let $e_{tail}$ be the final drop-off edge of the proposed subset $S$. We introduce a penalty for Euclidean or Network distance to the nearest active soft route $E_{soft}$.
  \[ C_{new} = C_{base}(v, S) \textbf{ } \mathbf{+} \textbf{ } \kappa \cdot \min_{e \in E_{soft}} \text{Distance}(e_{tail}, e) \]
  Where $\kappa$ is the penalty weight per meter.
  \end{block}
  
  \vspace{0.3cm}
  \textbf{Algorithmic Logic:}
  \begin{itemize}
      \item Acts as a mathematical "rubber band". The greedy algorithm will heavily penalize passenger combinations that pull the fleet away from the city's main demand arteries.
      \item Ensures vehicles are naturally pre-positioned for their next trip the moment they finish their current one.
  \end{itemize}
\end{frame}

\begin{frame}{Comparative Analysis \& Conclusion}
  \begin{table}[]
      \centering
      \small
      \begin{tabular}{|p{2.5cm}|p{4cm}|p{4cm}|}
      \hline
      \textbf{Approach} & \textbf{Primary Advantage} & \textbf{Computational Trade-off} \\ \hline
      \textbf{1. Adherence Reward} & Highly effective at maximizing fleet utilization along major corridors. & Requires calculating subset overlaps for every route permutation. \\ \hline
      \textbf{2. Spatial Detour} & Simplest to implement natively inside the TraCI/Dijkstra router algorithms. & Can lead to vehicles taking strange zigzag paths just to touch soft-route edges. \\ \hline
      \textbf{3. Terminal Cost} & Directly solves the "deadheading" problem by ensuring perfect drop-off placement. & Requires real-time nearest-neighbor calculations to the soft-route network. \\ \hline
      \end{tabular}
  \end{table}
  
  \vspace{0.3cm}
  \textbf{Recommendation for Next Steps:}
  \begin{itemize}
      \item \textbf{Approach 3 (Terminal Opportunity Cost)} offers the most elegant solution to the predictive rebalancing problem. It directly tackles the 86\% deadheading ratio by ensuring vehicles terminate trips inside active demand zones without breaking the natural flow of the simulation.
  \end{itemize}
\end{frame}

\end{document}
